% ===== Wrap-up (~4 min, then questions) =====
\partaccent{butter}\notopic

\begin{frame}{Այսօրվա 5 գաղափարը}
  \begin{enumerate}
    \item<1-> Ուղեղը կառուցված է ուրիշների \textbf{համար} ու ուրիշների \textbf{կողմից}. շփումն անհրաժեշտություն է
    \item<2-> Դեմքը խոսում է, նույնիսկ երբ բառերը լռում են. լիմբիկ համակարգը \textbf{ազնիվ} է
    \item<3-> Ոտքը դռան արանքում, դուռը երեսին, ցածր գնդակ. աշխատում են, որովհետև ուզում ենք, որ մեզ \textbf{սիրեն}
    \item<4-> Բաժանումը ցավում է, որովհետև զուգընկերը դարձել էր քո \textbf{մի մասը}։ Բայց ժամանակը բուժում է
    \item<5-> Խումբը կարող է ստիպել ժխտել սեփական աչքերը, հնազանդվել ու \textbf{մեղադրել զոհին}
  \end{enumerate}
\end{frame}

\begin{frame}{Փորձիր տանը}
  \begin{columns}[T]
    \begin{column}{0.31\textwidth}
      \begin{card}[colback=mint!30!cardwhite]
        \centering\bigemoji[1cm]{tv}\par\smallskip
        \raggedright 5 րոպե ֆիլմ նայիր \textbf{առանց ձայնի}։ Դեմքերից ու ժեստերից որքա՞ն ես հասկանում
      \end{card}
    \end{column}
    \begin{column}{0.31\textwidth}
      \begin{card}[colback=butter!45!cream]
        \centering\bigemoji[1cm]{carrot}\par\smallskip
        \raggedright Հաջորդ «մի փոքր խնդրանքի» կամ «հատուկ առաջարկի» ժամանակ գուշակիր, թե որ \textbf{հնարքն} է
      \end{card}
    \end{column}
    \begin{column}{0.31\textwidth}
      \begin{card}[colback=lavender!35!cardwhite]
        \centering\bigemoji[1cm]{speech}\par\smallskip
        \raggedright Երբ խմբային չաթում մեկն օգնություն է խնդրում, \textbf{մի սպասիր}, որ ուրիշը կպատասխանի
      \end{card}
    \end{column}
  \end{columns}
\end{frame}

\begin{frame}{Հաջորդ վեբինարը}
  \begin{columns}[c]
    \begin{column}{0.6\textwidth}
      \chip[lavender]{brain}{8. Երբ ուղեղը «փչանում» է}\par\medskip
      Շարքի վերջին վեբինարը։ Հոգեկան առողջություն. ինչ է իրականում դեպրեսիան ու ինչ թյուր պատկերացումներ կան դրա մասին, ինչպես է առաջանում «նյարդային պոռթկումը», ինչպես է ուղեղը կախվածություն ձևավորում, և որտեղից են գալիս հալյուցինացիաներն ու զառանցանքները։\par\medskip
      {\small\color{plum!70!cream}Մինչև հիմա տեսանք, թե ինչ տարօրինակ բաներ է անում ուղեղը, երբ \textbf{ամեն ինչ կարգին է}։ Հաջորդ անգամ՝ ինչ է լինում, երբ այն սկսում է խափանվել։}
    \end{column}
    \begin{column}{0.34\textwidth}
      \centering\mascot[2.8cm]{thinking}
    \end{column}
  \end{columns}
\end{frame}

\begin{frame}[plain]
  \begin{tikzpicture}[remember picture, overlay]
    \fill[lavender!60!cream] (current page.south west) rectangle (current page.north east);
    \fill[pattern={Dots[distance=4.4mm, radius=0.28mm]}, pattern color=plum!10!lavender!60!cream]
      (current page.south west) rectangle (current page.north east);
    \node at ([xshift=-3.2cm]current page.center) {\mascot[4.4cm]{happy}};
    \node[anchor=west, text=plum] at ([xshift=0.2cm, yshift=0.6cm]current page.center)
      {\titlefont\fontsize{40}{42}\selectfont Հարցեր};
    \node[anchor=west, text=plum, font=\mediumfont\large, align=left] at ([xshift=0.2cm, yshift=-1.0cm]current page.center)
      {Գրեք չաթում կամ\\ միացրեք խոսափողը \emoji{megaphone}};
    \node[anchor=west, text=plum!70!cream, font=\small] at ([xshift=0.2cm, yshift=-2.1cm]current page.center)
      {Դին Բըրնեթ, «The Idiot Brain», 2016};
  \end{tikzpicture}
\end{frame}
